\documentclass[11pt,a4paper]{article}
\usepackage[T1]{fontenc}
\usepackage[utf8]{inputenc}
\usepackage[polish]{babel}
\usepackage{lmodern}
\usepackage{geometry}
\usepackage{booktabs}
\usepackage{siunitx}
\usepackage{xcolor}
\usepackage{pgfplots}
\usepackage{hyperref}
\usepackage{caption}
\pgfplotsset{compat=1.18}
\geometry{margin=2.2cm}
\sisetup{group-separator={\,}}

\title{UniVerify: Wyniki Eksperymentów Kosztu Gazu}
\author{Nazar\\Projekt badawczy / praca magisterska}
\date{Wyniki z dnia 24 lutego 2026}

\begin{document}
\maketitle

\section{Streszczenie}
Przeprowadzono dwa etapy eksperymentów:
\begin{enumerate}
  \item optymalizacje kolejnych wersji kontraktu \texttt{A--F2},
  \item porównanie modelu \textit{legacy} (per-dyplom, \texttt{F2}) z modelem \textit{Merkle batch} (\texttt{DiplomaRegistry}).
\end{enumerate}
Wyniki pokazują, że A--F2 różnią się realnie kosztowo, a przejście na Merkle zmienia model kosztów: większy deployment, ale dużo niższy koszt \textit{na dyplom} przy większych batchach.

\section{Metodyka}
\begin{itemize}
  \item Środowisko: Foundry, Solidity \texttt{0.8.20}.
  \item Komendy pomiarowe:
\begin{verbatim}
cd contracts
forge test --match-path test/DiplomaRegistry*.t.sol --gas-report
forge test --match-path test/DiplomaRegistryMerkle.t.sol --gas-report
\end{verbatim}
  \item Dane: logi \texttt{forge} z testów A, B, C, D, E, F1, F2 oraz Merkle.
\end{itemize}

\section{Etap 1: wyniki A--F2}
\subsection{Suite-level (testy happy-path)}
\begin{table}[h]
\centering
\caption{Koszt testów \texttt{testGas\_addIssuer}, \texttt{testGas\_issue}, \texttt{testGas\_revoke}}
\begin{tabular}{lrrr}
\toprule
Wersja & addIssuer & issue & revoke \\
\midrule
A  & 57\,119 & 60\,721 & 93\,116 \\
B  & 57\,153 & 62\,890 & 114\,348 \\
C  & 57\,186 & 63\,012 & 114\,380 \\
D  & 55\,100 & 63\,037 & 114\,431 \\
E  & 55\,312 & 63\,047 & 114\,448 \\
F1 & 55\,100 & 60\,569 & 92\,868 \\
F2 & 55\,056 & 60\,531 & 92\,797 \\
\bottomrule
\end{tabular}
\end{table}

\noindent
\textbf{Uwaga metodologiczna:} \texttt{testGas\_revoke()} obejmuje także przygotowanie stanu (\texttt{issue} przed \texttt{revoke}), więc ta metryka nie jest czystym kosztem samej funkcji \texttt{revoke()}.

\begin{figure}[h]
\centering
\begin{tikzpicture}
\begin{axis}[
    width=0.95\textwidth,
    height=7cm,
    xlabel={Wersja kontraktu},
    ylabel={Gas},
    symbolic x coords={A,B,C,D,E,F1,F2},
    xtick=data,
    ymin=0,
    grid=both,
    legend style={at={(0.5,-0.2)},anchor=north,legend columns=2},
]
\addplot+[mark=*, very thick, color=blue!70!black] coordinates {(A,60721) (B,62890) (C,63012) (D,63037) (E,63047) (F1,60569) (F2,60531)};
\addlegendentry{testGas\_issue}
\addplot+[mark=square*, very thick, color=red!70!black] coordinates {(A,93116) (B,114348) (C,114380) (D,114431) (E,114448) (F1,92868) (F2,92797)};
\addlegendentry{testGas\_revoke}
\end{axis}
\end{tikzpicture}
\caption{Trend kosztów testów suite-level w wersjach A--F2.}
\end{figure}

\subsection{Function-level (czyste koszty funkcji, Avg)}
\begin{table}[h]
\centering
\caption{Porównanie wersji A--F2 na poziomie raportu funkcji}
\begin{tabular}{lrrrr}
\toprule
Wersja & Deployment & addIssuer & issue & revoke \\
\midrule
A  & 464\,211 & 45\,719 & 48\,458 & 31\,274 \\
B  & 469\,181 & 46\,963 & 50\,659 & 50\,369 \\
C  & 409\,497 & 45\,904 & 43\,367 & 38\,867 \\
D  & 403\,109 & 43\,366 & 43\,370 & 38\,870 \\
E  & 418\,029 & 43\,523 & 43\,377 & 38\,874 \\
F1 & 373\,484 & 43\,366 & 41\,607 & 28\,825 \\
F2 & 393\,596 & 43\,621 & 43\,084 & 29\,302 \\
\bottomrule
\end{tabular}
\end{table}

\subsection{Czym różnią się wersje A--F2 (z kodu)}
\begin{itemize}
  \item \textbf{A}: baza, \texttt{require(..., "STRING")}, rekord \texttt{\{issuer, issuedAt, revoked(bool)\}}.
  \item \textbf{B}: revokacja przez \texttt{revokedAt(uint64)} zamiast bool.
  \item \textbf{C}: przejście na \texttt{custom errors}.
  \item \textbf{D}: \texttt{owner} jako \texttt{immutable}.
  \item \textbf{E}: idempotentne \texttt{addIssuer/removeIssuer}, drobne porządki ścieżek.
  \item \textbf{F1}: usunięcie timestampów ze storage; rekord to \texttt{\{issuer, revoked\}}.
  \item \textbf{F2}: packed storage \texttt{mapping(bytes32 => uint256)} (issuer + bit revoked), plus \texttt{status()}.
\end{itemize}

\noindent
\textbf{Wniosek A--F2:} wersje realnie się różnią. Najsłabsza kosztowo jest \texttt{B}. Najlepszy kompromis legacy daje \texttt{F2}; \texttt{F1} ma najniższy deployment.

\section{Etap 2: F2 vs Merkle}
\subsection{Porównanie function-level (min/avg/max)}
\begin{table}[h]
\centering
\caption{Legacy \texttt{F2} -- raport funkcji}
\begin{tabular}{lrrr}
\toprule
Operacja & Min & Avg & Max \\
\midrule
\texttt{addIssuer} & 21\,585 & 43\,621 & 44\,855 \\
\texttt{issue} & 24\,030 & 43\,084 & 48\,200 \\
\texttt{removeIssuer} & 21\,607 & 23\,112 & 24\,947 \\
\texttt{revoke} & 26\,407 & 29\,302 & 31\,222 \\
\bottomrule
\end{tabular}
\end{table}

\begin{table}[h]
\centering
\caption{Nowy \texttt{DiplomaRegistry} (Merkle) -- raport funkcji}
\begin{tabular}{lrrr}
\toprule
Operacja & Min & Avg & Max \\
\midrule
\texttt{issueBatchRoot} & 24\,320 & 44\,433 & 50\,946 \\
\texttt{revokeFromBatch} & 30\,063 & 42\,624 & 55\,175 \\
\texttt{statusWithProof} & 5\,787 & 6\,071 & 6\,309 \\
\texttt{statusWithProofTrusted} & 8\,679 & 10\,193 & 10\,957 \\
\bottomrule
\end{tabular}
\end{table}

\subsection{Co zmienia Merkle (funkcjonalnie)}
\begin{itemize}
  \item Brak \texttt{issue(docHash)} per dyplom; zamiast tego \texttt{issueBatchRoot(batchId, root)}.
  \item Revokacja przez dowód Merkle: \texttt{revokeFromBatch(docHash, batchId, proof)}.
  \item Weryfikacja na proofach: \texttt{statusWithProof} / \texttt{statusWithProofTrusted}.
  \item Dodatkowo model uczelni: \texttt{issuerUniversityId}, status uczelni, \texttt{metadataHash}, \texttt{snapshotHash}.
\end{itemize}

\subsection{Amortyzacja kosztu publikacji}
Porównanie jednostkowe (wartości \texttt{Max}):
\[
\text{legacy: } \texttt{issue}_{max}=48\,200,\qquad
\text{Merkle: } \texttt{issueBatchRoot}_{max}=50\,946.
\]
\[
\text{gas\_per\_diploma}(N)=\frac{50\,946}{N}.
\]

\begin{figure}[h]
\centering
\begin{tikzpicture}
\begin{axis}[
    width=0.95\textwidth,
    height=7cm,
    xlabel={Rozmiar batcha \(N\)},
    ylabel={Gas / dyplom},
    xmode=log,
    log basis x={10},
    ymin=0,
    grid=both,
    legend style={at={(0.5,-0.2)},anchor=north,legend columns=2}
]
\addplot+[mark=*, very thick, color=green!50!black] coordinates {(1,50946) (10,5094.6) (100,509.46) (1000,50.946)};
\addlegendentry{Merkle: issueBatchRoot / N}
\addplot+[mark=square*, very thick, color=orange!90!black] coordinates {(1,48200) (10,48200) (100,48200) (1000,48200)};
\addlegendentry{Legacy: issue per dyplom}
\end{axis}
\end{tikzpicture}
\caption{Koszt jednostkowy publikacji dyplomu: legacy vs Merkle.}
\end{figure}

\section{Wnioski}
\begin{itemize}
  \item A--F2 to rzeczywiste optymalizacje kodu, nie kosmetyka.
  \item Największy regres kosztowy występuje w wersji \texttt{B}; najniższe koszty runtime osiągają \texttt{F1/F2}.
  \item Merkle to zmiana architektury, nie kolejna mikrooptymalizacja.
  \item Deployment Merkle jest wyższy (\texttt{1\,210\,329} vs \texttt{393\,596} dla F2), ale koszt operacyjny na dyplom spada wraz z \(N\).
  \item Dla \(N=100\): ok. 509 gas/dyplom (Merkle) vs 48\,200 gas/dyplom (legacy), czyli ok. \(94.8\times\) mniej.
\end{itemize}

\section{Ograniczenia}
\begin{itemize}
  \item Pomiary lokalne (Foundry), bez zmienności opłat sieciowych.
  \item Kolejny krok: eksperyment czułości \texttt{revokeFromBatch} na długość proof (\(O(\log N)\)).
\end{itemize}

\end{document}
